%% ============================================================
%%  IEEE Conference Proposal Paper
%%  Compile: pdflatex proposal.tex  (twice for cross-references)
%%  Class: IEEEtran (conference mode, two-column)
%% ============================================================

\documentclass[conference]{IEEEtran}
\IEEEoverridecommandlockouts

%% --- Packages ---
\usepackage{cite}
\usepackage{amsmath,amssymb,amsfonts}
\usepackage{algorithmic}
\usepackage{graphicx}
\usepackage{textcomp}
\usepackage{xcolor}
\usepackage{booktabs}
\usepackage{multirow}
\usepackage{url}
\usepackage{hyperref}
\usepackage{siunitx}
\sisetup{detect-all, per-mode=symbol}

\def\BibTeX{{\rm B\kern-.05em{\sc i\kern-.025em b}\kern-.08em
    T\kern-.1667em\lower.7ex\hbox{E}\kern-.125emX}}

%% ============================================================
\begin{document}

\title{Design of a Distributed LoRa-TDMA Power Telemetry Network
       with On-Node AC Load Classification and a Precision Analog
       Front-End for Residential Energy Monitoring}

\author{%
  \IEEEauthorblockN{L.\ A.\ Garol}
  \IEEEauthorblockA{\textit{College of Engineering} \\
    \textit{University of Mindanao} \\
    Davao City, Philippines \\
    l.garol.545854@umindanao.edu.ph}
}

\maketitle

%% ============================================================
\begin{abstract}
Residential electricity consumption in the Philippines remains poorly
characterised at the appliance level.\ Commercially available smart
meters provide aggregate kWh totals but offer no per-outlet load
identification or real-time waveform data, limiting both consumer
awareness and demand-side management programmes.\ This paper proposes
a multi-node wireless power telemetry system that addresses these gaps
through three integrated contributions.\ First, a precision analog
front-end based on the ADS131M02 24-bit simultaneous-sampling
delta-sigma ADC, a ZMPT101B galvanically isolated voltage transformer,
and a salvaged toroidal current transformer (CT) replaces the
commodity PZEM-004T v3 module.\ The new front-end provides
simultaneous voltage and current sampling at 4~kSPS with programmable
gain (G\,=\,1 to 128), full galvanic isolation from the AC mains, and
sufficient bandwidth for harmonic analysis up to the 15th order.\ Second, an
STM32F411 Cortex-M4F sensor node performs on-device harmonic feature
extraction via a 128-point FFT and classifies connected loads into six
residential categories---resistive, capacitive, single-phase motor,
fan, switched-mode power supply, and lighting---using a deterministic
decision tree with an upgradeable TFLite Micro multilayer perceptron
(MLP).\ Third, an ESP32-S3 gateway coordinates up to eight sensor
nodes through a frequency-hopping TDMA LoRa protocol with AES-128 CTR
packet encryption, hosts a web-based single-page application (SPA)
dashboard over Wi-Fi, and persists 120-point telemetry histories in
ferroelectric RAM (FRAM).\ A GFSK bulk transfer extension enables
over-the-air deployment of updated classifier models and calibration
coefficient blocks within the existing superframe structure.\ The
proposed system targets \(\le 0.5\%\) metering accuracy, \(\ge 85\%\)
classification accuracy across all six load categories, and
\(\ge 99\%\) uplink packet delivery ratio under residential indoor
conditions.
\end{abstract}

\begin{IEEEkeywords}
LoRa, TDMA, load classification, non-intrusive load monitoring,
ADS131M02, power quality, harmonic analysis, IoT, embedded machine
learning, smart metering
\end{IEEEkeywords}

%% ============================================================
\section{Introduction}
\label{sec:intro}

The residential sector accounts for a substantial and growing fraction
of national electricity consumption in the Philippines, yet the
granularity of available energy data remains coarse.\ Standard
electromechanical and electronic energy meters report only aggregate
kilowatt-hour totals, providing neither the per-outlet visibility
needed for occupant feedback nor the appliance-level signatures
required for automated demand management.\ Non-intrusive load
monitoring (NILM) techniques \cite{hart_nilm_1992} can, in principle,
disaggregate whole-house consumption into per-appliance contributions
from a single metering point, but whole-house NILM requires
significant signal processing resources and yields lower accuracy when
many loads operate concurrently \cite{zeifman_nilm_review_2011}.\ An
alternative approach---monitoring individual branch circuits or
outlets with networked smart plugs---trades installation complexity for
classification accuracy, since each measurement point sees at most a
few loads simultaneously.

Existing low-cost smart plug implementations typically rely on
purpose-built power metering ICs such as the PZEM-004T v3, which
integrates voltage sensing, current sensing, and RMS computation on a
single module.\ While sufficient for energy auditing, these modules
offer no access to the raw voltage or current waveform, precluding
harmonic analysis and thereby ruling out any feature-based load
classification.\ Replacing such a module with a discrete analog
front-end and a high-resolution simultaneous-sampling ADC restores
full waveform access at the cost of additional circuit complexity.

Long-range low-power wireless connectivity is a second challenge for
multi-point metering.\ Wi-Fi and Bluetooth are well-supported in the
Arduino and ESP-IDF ecosystems but impose high standby power and, for
Wi-Fi, require AP association with its associated latency and join
overhead.\ LoRa (Long Range), a chirp-spread-spectrum physical layer
operating in unlicensed ISM bands, achieves kilometre-scale range at
milliwatt transmit power and is increasingly used in sub-GHz IoT
deployments \cite{augustin_lora_sensors_2016}.\ A structured
medium-access control layer is required, however, to coordinate
multiple nodes without collision; frequency-division and
time-division multiple access (TDMA) schemes have been demonstrated to
be effective for deterministic, low-latency sensor networks
\cite{raza_lpwan_survey_2017}.

Security is a practical concern for metering infrastructure.\ An
unencrypted LoRa broadcast reveals occupancy patterns, appliance usage
schedules, and potentially security-relevant information (vacation
periods, nighttime routines).\ AES-128 in counter (CTR) mode provides
a lightweight, nonce-based stream cipher that is compatible with
embedded C++ without dynamic memory allocation and is recommended by
NIST \cite{nist_sp800_38a}.

This paper proposes a system that integrates all three of these
requirements: precision analog measurement, structured wireless MAC,
and load classification---within a hardware/firmware architecture that
is practical to build and deploy in a Philippine residential setting.

\subsection{Objectives}
\label{ssec:objectives}

The specific objectives of this research are:

\begin{enumerate}
  \item Design and prototype a precision analog front-end using the
        ADS131M02 24-bit simultaneous-sampling ADC with a ZMPT101B
        isolated voltage transformer and a toroidal CT, targeting
        \(\le 0.5\%\) metering accuracy and full galvanic isolation from
        the AC mains.

  \item Implement a TDMA LoRa communication protocol with
        frequency-hopping channel access, AES-128 CTR packet
        encryption, and a GFSK bulk transfer extension for
        over-the-air model and calibration deployment.

  \item Develop an on-node harmonic feature extraction pipeline using
        a 128-point FFT and derive a 10-dimensional feature vector
        for load identification without transmitting raw waveform data.

  \item Implement and evaluate a two-stage on-node load classifier: a
        deterministic decision tree for immediate deployment and a
        compact TFLite Micro MLP trained on collected field data,
        targeting \(\ge 85\%\) accuracy across six residential load
        categories.

  \item Develop an ESP32-S3 gateway with a web-based SPA dashboard
        for real-time visualization, remote relay control, persistent
        energy history via FRAM, and over-the-air classifier
        management.

  \item Validate the integrated system against reference measurements,
        characterizing metering accuracy, classification accuracy, and
        uplink packet delivery ratio under realistic residential
        conditions.
\end{enumerate}

\subsection{Scope and Limitations}
\label{ssec:scope}

The system targets single-phase 220~V / 60~Hz residential loads of up
to 30~A per monitored outlet.\ Gateway-to-internet connectivity
(cloud backhaul) is out of scope; the dashboard is self-hosted on the
gateway node.\ Load disaggregation of mixed-load outlets is not
addressed.\ The analogue front-end is designed for indoor perfboard
prototype construction; PCB layout and EMC compliance are deferred.

%% ============================================================
\section{Related Work}
\label{sec:related}

\subsection{Non-Intrusive Load Monitoring and Load Classification}

Hart's foundational work \cite{hart_nilm_1992} established the
framework for identifying appliances from changes in steady-state
current and power at a single metering point.\ The approach works well
for stepwise loads but struggles with variable-speed drives and
modern switched-mode supplies.\ Subsequent work extended feature sets
to include harmonic content \cite{wichakool_harmonic_loads_2009},
transient event detection, and V-I trajectory analysis
\cite{lam_vI_trajectory_2007}.\ Zeifman and Roth
\cite{zeifman_nilm_review_2011} survey the state of the art and
identify harmonic-ratio features as among the most discriminative for
residential load types.

Data-driven approaches have gained traction as labelled datasets
became available.\ Kolter and Johnson \cite{kolter_redd_2011}
released the REDD dataset and demonstrated that hidden Markov models
can exceed the accuracy of threshold-based disaggregation.\ Kelly and
Knottenbelt \cite{kelly_neural_nilm_2015} applied recurrent neural
networks directly to raw power traces, achieving state-of-the-art
disaggregation on the UK-DALE dataset.\ For per-outlet classification
(as opposed to whole-house disaggregation), the classification problem
is considerably simpler because each metering point typically
serves one or a few loads; hand-crafted feature-based classifiers
have been shown to reach \(>90\%\) accuracy in this setting
\cite{parson_unsupervised_nilm_2012}.

The proposed system adopts a hybrid approach: a deterministic
decision tree provides immediate classification without training data,
while a compact MLP---deployed via OTA after supervised data
collection---improves accuracy on atypical or region-specific loads.

\subsection{Precision Power Measurement with Sigma-Delta ADCs}

High-resolution simultaneous-sampling ADCs have become the preferred
architecture for precision power quality instruments.\ The ADS131M02
(Texas Instruments) \cite{ti_ads131m02_ds} is a 24-bit, dual-channel,
simultaneously sampling delta-sigma ADC with an on-chip PGA (gain
1--128) and a sinc3 decimation filter, delivering 18--20 effective
bits at output data rates of 1--32~kSPS.\ The simultaneous sampling
architecture eliminates the phase skew inherent in multiplexed ADCs,
which is critical for accurate reactive-power and power-factor
measurement.

Several precision metering IC designs reported in the literature use
this class of ADC.\ Adamo et al.\ \cite{adamo_power_quality_2011}
describe a power quality monitor based on a 16-bit simultaneous ADC,
demonstrating that waveform-based harmonic analysis is practical in
embedded systems.\ The ZMPT101B miniature current-mode voltage
transformer provides 4~kV galvanic isolation from the AC mains with
less than 20 arcminutes of phase error at the rated operating current,
making it a cost-effective alternative to resistive dividers for safe
isolated voltage sensing \cite{zmpt101b_datasheet}.

\subsection{LoRa and LPWAN for IoT Sensing}

LoRa modulation, developed by Semtech \cite{semtech_sx1278_ds}, uses
chirp spread spectrum (CSS) over configurable spreading factors (SF6
to SF12), trading data rate for link margin.\ At SF6 the air-time
for a 32-byte packet is approximately 10--15~ms with \(\sim 10\)~dB
less processing gain than SF12; this is appropriate for TDMA
networks where short slot durations are required.\ Augustin et al.\
\cite{augustin_lora_sensors_2016} characterise LoRa coverage in urban
and indoor environments, confirming reliable indoor coverage at 20--30~m
through multiple concrete walls at SF6 with 10~dBm transmit power.

TDMA-based MAC layers for LoRa have been proposed to overcome the
collision-prone LoRaWAN ALOHA-like access scheme.\ Bor et al.\
\cite{bor_lorawan_scale_2016} show that LoRaWAN does not scale beyond
a few hundred devices in dense deployments, motivating structured MAC
protocols for high-density or latency-sensitive applications.\ Frequency
hopping across multiple channels increases resilience against
narrowband interference and partial-band fading \cite{raza_lpwan_survey_2017}.

\subsection{Embedded Machine Learning for IoT}

The TFLite Micro runtime \cite{david_tflite_micro_2021} enables
inference from quantised TensorFlow Lite models on microcontrollers
with as little as 16~kB of RAM.\ Warden and Situnayake
\cite{warden_tinyml_2019} survey the domain and demonstrate that small
(hundreds of parameters) multilayer perceptrons achieve adequate
accuracy for many IoT classification tasks while executing in under
1~ms on Cortex-M4 processors.\ The ARM CMSIS-DSP library
\cite{arm_cmsis_dsp} provides optimised FFT and vector operations for
Cortex-M processors, enabling efficient feature extraction without
hand-tuned assembly.

\subsection{Security in LoRa IoT Networks}

LoRaWAN mandates AES-128 for join-request authentication and
payload encryption, demonstrating that AES is a practical security
primitive for constrained IoT devices \cite{lorawan_spec}.\ NIST
SP~800-38A \cite{nist_sp800_38a} defines CTR mode, which converts a
block cipher into a stream cipher suitable for variable-length packets
and allows parallelised encryption without IV re-use if the nonce is
unique per packet.\ Hardware RFID-based key provisioning, as studied
by Hancke and Kuhn \cite{hancke_rfid_security_2005}, provides a
physical-world complement to over-the-air key distribution.

%% ============================================================
\section{System Architecture}
\label{sec:arch}

Fig.\ \ref{fig:arch} illustrates the proposed system.\ Two hardware
roles share a single firmware codebase compiled under different
PlatformIO build environments.

\begin{figure}[t]
  \centering
  %% Replace with actual figure:
  %% \includegraphics[width=\columnwidth]{figs/system_arch.pdf}
  \framebox{\parbox{\columnwidth}{\centering
    \vspace{2cm}
    [System Architecture Diagram]\\
    Gateway (ESP32-S3) $\leftarrow$ LoRa TDMA $\rightarrow$ Nodes $\times$8
    (STM32F411 + ADS131M02)\\
    Gateway $\leftarrow$ Wi-Fi $\rightarrow$ Web Dashboard
    \vspace{2cm}
  }}
  \caption{Proposed system architecture.\ The ESP32-S3 gateway acts
    as the TDMA master and Wi-Fi access point.\ Up to eight
    STM32F411 sensor nodes report simultaneously sampled voltage and
    current waveforms and classified load types every 3~s superframe.}
  \label{fig:arch}
\end{figure}

\textbf{Gateway (ESP32-S3):} The gateway is built on the Espressif
ESP32-S3 SoC (Xtensa LX7 dual-core at 240~MHz, 512~kB SRAM).\ It
acts as the TDMA master, transmitting a beacon at the start of each
3~s superframe and allocating uplink/downlink slots to registered
nodes.\ A Wi-Fi STA/AP interface serves the management dashboard.\ An
MB85RC256V 32~kB FRAM chip (I\textsuperscript{2}C, 0x50) persists
per-node energy accumulators and 120-point telemetry histories across
reboots.\ An optional PN532 NFC reader supports hardware-based AES key
provisioning in encrypted builds.

\textbf{Sensor Nodes (STM32F411):} Each sensor node uses the
STMicroelectronics STM32F411 ``Black Pill'' (ARM Cortex-M4F at
100~MHz, 128~kB SRAM).\ The node samples the AC mains via the custom
analog front-end (Section~\ref{ssec:afe}), runs the on-node load
classifier (Section~\ref{ssec:classifier}), controls a relay for
switched load management, and transmits a 32-byte telemetry packet in
its assigned TDMA uplink slot.

\textbf{Radio Subsystem:} Both gateway and nodes use the Semtech
SX1278 LoRa transceiver connected via SPI.\ The gateway uses the
standard ESP32-S3 FSPI peripheral; the nodes use STM32 SPI1 for the
ADS131M02 ADC and SPI2 for the SX1278.\ The RadioLib library
\cite{radiolib} abstracts low-level transceiver control.

%% ============================================================
\section{Methodology}
\label{sec:method}

\subsection{Analog Front-End Design}
\label{ssec:afe}

\subsubsection{Current Channel}

The current channel uses the CT bundled with the PZEM-004T module---a
solid-core toroidal transformer with a 1:1000 turns ratio and
100~A primary rating.\ The CT secondary drives a split burden resistor
(\(2 \times 5\,\Omega\), 0.1\%, 15~ppm/\textdegree C metal film),
whose center tap is connected to AGND.\ This balanced configuration
eliminates even-order distortion in the ADC input stage.

At the 30~A target load, the peak burden voltage is:
\begin{equation}
  V_{b,\text{pk}} = \frac{I_{\max}\sqrt{2}}{n} \cdot R_b
                  = \frac{30\sqrt{2}}{1000} \times 10 = 0.424\,\text{V},
  \label{eq:burden_v}
\end{equation}
corresponding to 70.7\% of the ADS131M02 full-scale range at PGA
gain \(G = 2\) (\(\pm 0.6\)~V).\ The clipping threshold at \(G = 2\)
is 42.4~A RMS, accommodating steady-state residential loads; motor
inrush transients up to 84.9~A RMS are captured by switching to
\(G = 1\) in firmware.\ The noise-limited current resolution at
4~kSPS (\(G = 2\), \(e_n \approx 3.5\,\mu\text{V}\)) is:
\begin{equation}
  \Delta I_{p} = \frac{e_n}{R_b} \cdot n
               = \frac{3.5 \times 10^{-6}}{10} \times 1000
               = 0.35\,\text{mA},
  \label{eq:i_resolution}
\end{equation}
yielding a dynamic range of approximately 60\,000:1, far exceeding the
5\,000:1 design target.

\subsubsection{Voltage Channel}

The ZMPT101B miniature current-mode voltage transformer (1:1 ratio,
4~kV isolation) senses the mains voltage.\ Two 56~k\(\Omega\) series
resistors (\(R_L = 112\,\text{k}\Omega\) total) limit the primary
operating current to 2~mA at 220~V RMS.\ The secondary current flows
through a split sampling resistor (\(2 \times 160\,\Omega\), center
tap to AGND), producing a peak differential output of 0.917~V at
220~V RMS---76.4\% of the ADS131M02 full-scale at \(G = 1\).\ The
ZMPT101B introduces a phase lag of \(\le 0.33^\circ\), characterised in
\cite{zmpt101b_datasheet}.

\subsubsection{Anti-Aliasing Filter and Phase Calibration}

Both channels share an identical differential RC anti-aliasing filter:
\(R_{\text{filt}} = 100\,\Omega\) on each line, \(C_{\text{diff}} =
680\,\text{nF}\) C0G ceramic across the differential pair, giving:
\begin{equation}
  f_c = \frac{1}{2\pi R_{\text{filt}} C_{\text{diff}}}
      = \frac{1}{2\pi \times 100 \times 680 \times 10^{-9}}
      = 2341\,\text{Hz}.
  \label{eq:filter_fc}
\end{equation}
Attenuation at the fundamental (60~Hz) is 0.003~dB; at the 15th
harmonic (900~Hz) it is 0.59~dB---a systematic, known, and correctable
gain error.\ Because both channels use identical filter components,
the differential phase error contributed by the anti-aliasing filter
is exactly zero \cite{ti_ads131m02_ds}; the ADS131M02's simultaneous
sampling eliminates ADC-induced phase skew.\ Residual differential
phase error from the ZMPT101B and CT transformers is corrected via the
ADS131M02 on-chip PHASE register at a resolution of:
\begin{equation}
  \Delta\phi_{\text{step}} = 360^\circ \times f_{\text{mains}}
                            \times \frac{1}{f_{\text{CLKIN}}}
                           = 360 \times 60 \times \frac{1}{8.192 \times 10^6}
                           \approx 0.00264^\circ.
  \label{eq:phase_step}
\end{equation}

\subsubsection{PGA Gain-Ranging}

Four gain states are defined in firmware (Table~\ref{tab:pga}).\ At
each 60~Hz cycle boundary, the firmware checks whether the peak sample
amplitude exceeds 90\% of the current FSR (switch to lower gain) or
falls below 15\% (switch to higher gain).\ Gain transitions are
applied via SPI register write and take effect on the next ADS131M02
conversion.\ The strategy maintains high ADC utilisation across a
5000:1 current range without increasing burden resistance.

\begin{table}[t]
  \caption{PGA Gain-Ranging States -- Current Channel}
  \label{tab:pga}
  \centering
  \begin{tabular}{cccc}
    \toprule
    Gain & FSR ($\pm$V) & Clip (A RMS) & Noise Floor (mA) \\
    \midrule
    1  & 1.20 & 84.9 & 0.53 \\
    2  & 0.60 & 42.4 & 0.35 \\
    4  & 0.30 & 21.2 & 0.20 \\
    16 & 0.075 & 5.3 & 0.095 \\
    \bottomrule
  \end{tabular}
\end{table}

\subsubsection{Firmware Integration on STM32F411}

The ADS131M02 DRDY pin is connected to an EXTI interrupt (priority 0,
highest).\ The ISR launches a DMA-driven SPI1 read of the 80-bit frame
(status word, 24-bit Ch0, 24-bit Ch1, CRC-16).\ A DMA completion
callback copies the two signed 24-bit samples into circular buffers,
waking the \texttt{adcTask} (FreeRTOS, STM32duino) for per-sample
accumulation.\ At 4~kSPS the total CPU load from ADC acquisition and
per-sample arithmetic is approximately 2.2\%, leaving \(>97\%\) idle
for the TDMA radio task and classifier.

\subsection{TDMA Protocol and Wireless Communication}
\label{ssec:tdma}

\subsubsection{Superframe Structure}

The network organises time into 3000~ms superframes (Fig.\
\ref{fig:superframe}).\ Each superframe contains:
\begin{itemize}
  \item A 40~ms beacon (gateway \(\to\) all nodes) on Ch~0 at
        SF6, 125~kHz BW.
  \item Eight 165~ms data slots (DL 50~ms + UL 100~ms + 15~ms guard)
        for eight registered nodes.
  \item A 60~ms contention window (uplink) and a 40~ms contention
        window (downlink) for unregistered nodes to join.
  \item A 1520~ms idle window used for gateway maintenance and,
        when a bulk transfer is active, a GFSK burst phase.
  \item A 20~ms end guard before the next beacon.
\end{itemize}

\begin{figure}[t]
  \centering
  \framebox{\parbox{\columnwidth}{\centering
    \vspace{1.5cm}
    [Superframe timing diagram:\\
    Beacon (40ms) | Slots 1-8 (8x165ms=1320ms) |\\
    CW-UL (60ms) | CW-DL (40ms) | Idle (1520ms) | Guard (20ms)]
    \vspace{1.5cm}
  }}
  \caption{TDMA superframe structure.\ The 1520~ms idle window is
    normally quiescent but is repurposed for GFSK bulk data transfers
    when a bulk session is active.}
  \label{fig:superframe}
\end{figure}

Downlink-first slot ordering gives the gateway the opportunity to
deliver commands (relay toggle, schedule, threshold) immediately
before the node's uplink.\ The node's uplink slot timing is derived
from the beacon receive timestamp:
\begin{equation}
  t_{\text{tx}} = t_{\text{bcn}} + (T_{\text{bcn}} - T_{\text{air}})
                + (k-1) \cdot T_{\text{slot}} + T_{\text{DL}},
  \label{eq:slot_time}
\end{equation}
where \(T_{\text{bcn}} = 40\,\text{ms}\), \(T_{\text{air}} =
18\,\text{ms}\) (on-air time of the 8-byte beacon at SF6),
\(T_{\text{slot}} = 165\,\text{ms}\), \(T_{\text{DL}} = 50\,\text{ms}\),
and \(k\) is the node's assigned slot index (1--8).

\subsubsection{Frequency Hopping}

The SX1278 is configured with eight 200~kHz-spaced channels in the
433~MHz ISM band (Ch~0 = 433.050~MHz, Ch~7 = 434.450~MHz).\ The
active channel for each slot within a superframe is:
\begin{equation}
  \text{ch} = (\mathit{sfCount} \times 7 + k) \bmod 8,
  \label{eq:freqhop}
\end{equation}
where \(\mathit{sfCount}\) is the superframe counter and the coprime
multiplier 7 ensures that every slot visits all eight channels across
eight consecutive superframes.\ Beacon and contention windows use
Ch~0.\ The hopping sequence is deterministic and requires no explicit
signalling.

\subsubsection{Node Lifecycle and Network Epoch}

Unregistered nodes transmit a \texttt{JoinRequest} (0xA0, 4~bytes)
during the contention window.\ The gateway allocates the first
available slot and replies with a \texttt{JoinAck} (0xA1, 4~bytes)
containing the assigned slot index and the current network epoch.\ A
uint8\_t epoch field in the beacon increments on every node eviction
(which occurs only when all 8 slots are occupied and a node has missed
8 consecutive superframes~\(\approx 24\)~s).\ A registered node that
detects an epoch mismatch immediately re-contends rather than
transmitting in a potentially reassigned slot.

\subsubsection{Packet Encryption}

When the \texttt{PKT\_ENCRYPTION} build flag is active, all
uplink and downlink packets (excluding the plaintext \texttt{sfCount}
bytes of the beacon) are encrypted in-place using AES-128 in CTR mode
\cite{nist_sp800_38a}.\ The 16-byte nonce is constructed as:
\begin{equation}
  \text{nonce} = [\,\mathit{sfCount}_\text{lo} \mid
                    \mathit{sfCount}_\text{hi} \mid
                    \mathit{slotId} \mid
                    \mathit{dir} \mid
                    \mathbf{0}_{12}\,],
  \label{eq:nonce}
\end{equation}
where \(\mathit{dir}\) is a per-packet-type direction discriminator
byte.\ This guarantees a unique keystream block per (superframe,
slot, direction) triple.\ The AES key is auto-generated by the gateway
on first boot and provisioned to nodes via a PN532 I\textsuperscript{2}C
NFC reader using MIFARE Classic sector~1 (Key~A 0xFF$\times$6) as
described in \cite{hancke_rfid_security_2005}.

\subsection{GFSK Bulk Transfer Extension}
\label{ssec:gfsk}

Several subsystems require payloads of 28--4096~bytes that cannot be
delivered efficiently through the 7-byte downlink slot window: ADC
calibration coefficient blocks (24~bytes), classifier threshold sets
(28~bytes), quantised TFLite Micro MLP models (\(\approx 500\)~bytes),
and training-data upload batches (2--4~kB).\ Delivering a 500-byte
payload through the 7-byte DL window would require 72 superframes
(216~s), which is impractical.

The GFSK bulk transfer extension repurposes the 1520~ms idle window
for a high-throughput GFSK burst, while leaving all eight TDMA data
slots entirely undisturbed (Fig.\ \ref{fig:bulk}).\ The protocol is
application-agnostic: it delivers an opaque byte buffer between the
gateway and one target node per superframe.

\begin{figure}[t]
  \centering
  \framebox{\parbox{\columnwidth}{\centering
    \vspace{1.5cm}
    [Bulk transfer timing diagram:\\
    ...CW-DL (40ms) | Maint (20ms) | LoRa->GFSK | GFSK Burst |
    GFSK->LoRa | Idle | Guard]
    \vspace{1.5cm}
  }}
  \caption{GFSK bulk transfer phase within the idle window.\ TDMA data
    slots and contention windows are completely unaffected.}
  \label{fig:bulk}
\end{figure}

\subsubsection{Physical Layer}

The SX1278 supports FSK/GFSK packet mode alongside LoRa through the
\texttt{LongRangeMode} bit in \texttt{RegOpMode}.\ Mode switching
requires a transition through Sleep mode and takes approximately
0.5~ms.\ GFSK is configured at 100~kbps bit rate, 50~kHz frequency
deviation (modulation index \(\beta = 1.0\)), 125~kHz single-sided RX
bandwidth, BT~=~0.5 Gaussian filtering, and CRC-16/CCITT hardware
error detection.\ At 100~kbps, a worst-case 4096-byte transfer
(21 fragments of 200~bytes each) completes in approximately 462~ms,
well within the 1476~ms usable bulk window.

The indoor link budget at 433~MHz over 20~m through two concrete walls
is estimated at 73~dB path loss.\ With 10~dBm transmit power and
\(-104\)~dBm sensitivity at 100~kbps, the margin is 41~dB---adequate
for body shadowing and transient fading \cite{augustin_lora_sensors_2016}.

\subsubsection{Fragment Protocol}

The bulk payload is segmented into 200-byte fragments, each
encapsulated in a \texttt{BulkFragHeader} (pktType 0xB0, sequence
number, fragment count, payload length) followed by the payload and a
hardware CRC-16.\ A stop-and-wait ARQ scheme with a 50~ms ACK timeout
and up to 3 retransmissions per fragment provides reliable delivery.
An end-to-end CRC-32 over the complete reassembled payload detects
any reassembly error that per-fragment CRC would miss.

In encrypted builds, each fragment uses an extended nonce
incorporating the fragment sequence number in byte 4, guaranteeing
unique keystreams across fragments within a transfer and across
retried transfers in different superframes:
\begin{equation}
  \text{nonce}_{\text{bulk}} = [\,\mathit{sfCount}_\text{lo} \mid
                                   \mathit{sfCount}_\text{hi} \mid
                                   \mathit{slotId} \mid
                                   \mathit{dir} \mid
                                   \mathit{seqNum} \mid
                                   \mathbf{0}_{11}\,].
  \label{eq:bulk_nonce}
\end{equation}

The gateway signals the start of a bulk session by inserting a
6-byte \texttt{BulkGrantPacket} (pktType 0x08) into the target node's
regular LoRa downlink slot.\ The node acknowledges readiness by
setting bit~7 of the \texttt{statusFlags} byte in its subsequent
uplink telemetry packet.\ If the gateway observes this flag, both
sides simultaneously switch to GFSK at the computed bulk phase start
time; otherwise, the grant is retried in the next superframe.

\subsection{Feature Extraction and Load Classification}
\label{ssec:classifier}

\subsubsection{Feature Extraction Pipeline}

Every 2~s (approximately 120 mains cycles at 60~Hz), the
\texttt{adcTask} on the STM32F411 applies a 128-point FFT to the most
recent current waveform samples.\ With \(f_s = 4000\,\text{SPS}\) and
\(N = 128\), the FFT frequency resolution is:
\begin{equation}
  \Delta f = \frac{f_s}{N} = \frac{4000}{128} = 31.25\,\text{Hz}.
  \label{eq:fft_res}
\end{equation}
The fundamental (60~Hz) maps to bin \(k_1 = 2\); odd harmonics up to
the 15th (900~Hz) map to bins 2, 6, 10, 13, 17, 21, 25, 29.\ A Hann
window reduces spectral leakage; since only magnitude ratios are used
for classification, the 6~dB main-lobe broadening is acceptable.\ The
anti-aliasing filter's known attenuation at each harmonic frequency is
compensated by the inverse filter magnitude:
\begin{equation}
  |I_{h,\text{corr}}| = |I_h| \times \sqrt{1 + \left(\frac{h \cdot 60}{f_c}\right)^2}.
  \label{eq:filt_correction}
\end{equation}

From the corrected harmonic magnitudes, a 10-dimensional feature vector
\(\mathbf{x} = [f_0, \ldots, f_9]\) is computed (Table~\ref{tab:features}).
All features are dimensionless ratios or normalised quantities, making
the classifier independent of absolute load magnitude.\ The total feature
vector occupies 40~bytes.

\begin{table}[t]
  \caption{Feature Vector Definition}
  \label{tab:features}
  \centering
  \small
  \begin{tabular}{clp{4.0cm}}
    \toprule
    Index & Symbol & Definition \\
    \midrule
    0 & $f_\text{PF}$   & Displacement power factor \\
    1 & $f_\text{Qsgn}$ & Sign of reactive power ($\pm1$, 0) \\
    2 & $f_\text{THD}$  & $\text{THD}_I = \sqrt{\sum |I_h|^2}\,/\,|I_1|$ \\
    3 & $f_\text{H3}$   & $|I_3|\,/\,|I_1|$ \\
    4 & $f_\text{H5}$   & $|I_5|\,/\,|I_1|$ \\
    5 & $f_\text{H7}$   & $|I_7|\,/\,|I_1|$ \\
    6 & $f_\text{CF}$   & Crest factor $I_\text{pk}/I_\text{rms}$ \\
    7 & $f_\text{OE}$   & $(|I_2|+|I_4|)\,/\,|I_1|$ (asymmetry) \\
    8 & $f_\text{H53}$  & $|I_5|\,/\,|I_3|$ (shape discriminator) \\
    9 & $f_\text{Prms}$ & $\log_2(I_\text{rms}/0.05)$, clamped $[0,10]$ \\
    \bottomrule
  \end{tabular}
\end{table}

\subsubsection{Six-Class Load Taxonomy}

Six load categories are defined to cover the dominant residential load
types in a 220~V / 60~Hz Philippine residential installation
(Table~\ref{tab:classes}).\ The taxonomy is designed so that each class
has at least one strongly discriminative feature: SMPS loads are
identified by their high crest factor (\(> 2.2\)); lighting (CFL/LED)
by high H3/H1 ratio (\(> 0.35\)); motor loads by lagging power factor
with low THD; fans by deep lagging power factor; capacitive loads by
negative reactive power; and resistive loads by near-unity power factor
and low THD.

\begin{table}[t]
  \caption{Six-Class Load Taxonomy}
  \label{tab:classes}
  \centering
  \small
  \begin{tabular}{clp{3.4cm}}
    \toprule
    ID & Class & Typical Loads \\
    \midrule
    0 & RESISTIVE  & Incandescent, flat iron, water heater \\
    1 & CAPACITIVE & PFC-corrected SMPS, UPS on battery \\
    2 & MOTOR\_1PH & Aircon compressor, refrigerator \\
    3 & FAN        & Ceiling fan, desk fan, exhaust fan \\
    4 & SMPS       & Laptop/phone charger, desktop PSU \\
    5 & LIGHTING   & CFL, LED lamp (non-PFC), fluorescent \\
    \bottomrule
  \end{tabular}
\end{table}

\subsubsection{Baseline: Deterministic Decision Tree}

The primary classifier is a hand-crafted decision tree that requires
no labelled training data, executes in under 1~\(\mu\)s on the
STM32F411, and is fully interpretable.\ The tree follows a
``peel off the easy cases first'' strategy (Fig.~\ref{fig:dtree}):

\begin{enumerate}
  \item \textbf{No-load gate} (\(I_\text{rms} < 50\,\text{mA}\)):
        below the noise floor; classify as RESISTIVE (idle).
  \item \textbf{THD split} (\(\text{THD}_I < 0.15\)):
        separates approximately sinusoidal loads (motor, fan,
        resistive, capacitive) from distorted loads (SMPS, lighting).
  \item \textbf{Low-THD branch}:
        PF \(> 0.95\) \(\Rightarrow\) RESISTIVE;
        \(Q_\text{sgn} < 0\) \(\Rightarrow\) CAPACITIVE;
        PF \(> 0.72\) \(\Rightarrow\) MOTOR\_1PH; else FAN.
  \item \textbf{High-THD branch}:
        CF \(> 2.2\) \(\Rightarrow\) SMPS;
        H3/H1 \(> 0.35\) \(\Rightarrow\) LIGHTING;
        residual: apply PF and \(Q_\text{sgn}\) as in step~3.
\end{enumerate}

\begin{figure}[t]
  \centering
  \framebox{\parbox{\columnwidth}{\centering
    \vspace{2cm}
    [Decision Tree Diagram:\\
    Root: Irms<50mA? -> RESISTIVE\\
    THD<0.15? -> Low-THD branch (PF, Qsign) / High-THD branch (CF, H3)
    ]
    \vspace{2cm}
  }}
  \caption{Simplified decision tree structure.\ Each threshold is
    stored in a \texttt{ClassifierThresholds} struct and can be updated
    via OTA bulk transfer from the gateway.}
  \label{fig:dtree}
\end{figure}

A confidence score (0--7, packed into 3 bits) is computed as the
normalised distance from the crossed threshold.\ A 5-sample temporal
smoother (confidence-weighted majority vote over 10~s of history)
suppresses transient misclassifications during motor startup.\ The
complete classification result, including confidence and transient
flag, is packed into a single byte within the 32-byte
\texttt{TelemetryPacket}.

\subsubsection{Upgrade Path: TFLite Micro MLP}

After supervised training data are collected (Phase~1, Section~\ref{sec:validation}),
the decision tree is replaced by a compact 3-layer MLP (10 inputs,
16 hidden (ReLU), 12 hidden (ReLU), 6 outputs (softmax)) with 458
parameters.\ Post-training 8-bit quantisation reduces the model to
approximately 500~bytes, deployable over the GFSK bulk channel.\ The
TFLite Micro runtime requires a 4~kB tensor arena; inference takes
approximately 2~\(\mu\)s on the STM32F411 Cortex-M4F using
CMSIS-DSP-accelerated matrix operations \cite{arm_cmsis_dsp,
warden_tinyml_2019}.\ A fallback mechanism reverts to the decision
tree when the MLP softmax probability is below 0.5.

\subsection{Gateway Dashboard and Data Management}
\label{ssec:gateway}

The ESP32-S3 gateway hosts a LittleFS single-page application (SPA)
comprising HTML, CSS, JavaScript, and Chart.js, served via the
ESPAsyncWebServer library.\ The SPA communicates with the gateway
through a REST API and a WebSocket endpoint.\ Table~\ref{tab:api}
summarises the key REST endpoints.

\begin{table}[t]
  \caption{Gateway REST and WebSocket API (Selected)}
  \label{tab:api}
  \centering
  \small
  \begin{tabular}{lp{5.0cm}}
    \toprule
    Endpoint & Function \\
    \midrule
    \texttt{GET /api/nodes}           & All registered node states \\
    \texttt{GET /api/node/\{id\}/history} & 120-point V/I/P/class history \\
    \texttt{POST /api/bulk}           & Enqueue GFSK bulk DL/UL transfer \\
    \texttt{GET /api/features}        & Active build-flag feature set \\
    \texttt{POST /api/login}          & Issue session token \\
    WS: \texttt{telemetry}            & Real-time node update push \\
    WS: \texttt{relay\_manual}        & Toggle relay on/off \\
    WS: \texttt{relay\_schedule}      & Set timed relay schedule \\
    WS: \texttt{set\_threshold}       & Update power-off threshold \\
    \bottomrule
  \end{tabular}
\end{table}

Per-node energy accumulators, 120-point telemetry histories
(encoded as \texttt{HistoryPoint} structs at 16~bytes each), and node
labels are persisted in the MB85RC256V FRAM.\ The memory map allocates
1968~bytes per node slot (8~slots), occupying 15.7~kB of the 32~kB
chip.\ FRAM writes are deferred through a dirty-count mechanism
(threshold 10) and dispatched to a dedicated Core~0 FRAM task via a
FreeRTOS queue, ensuring that I\textsuperscript{2}C writes never block
the Core~1 TDMA timing task.\ Cross-slot UID lookup allows the
gateway to restore a node's energy history and label after a reboot
even if the node is assigned a different slot index.

The dashboard password is stored in ESP-IDF NVS (namespace
\texttt{wifi-cfg}, key \texttt{dashpass}).\ When a password is set,
all REST and WebSocket endpoints require a session token obtained from
\texttt{POST /api/login}.\ The system also supports optional AP
password protection, NTP-synchronized timestamps, and static IP
configuration.

\subsection{System Integration}
\label{ssec:integration}

The firmware is compiled from a single codebase under ten PlatformIO
environments differentiating gateway vs.\ node role, encrypted vs.\
plaintext, hardware PZEM vs.\ simulated sensor, and ESP32-dev vs.\
ESP32-S3 vs.\ STM32F411 target.\ Shared packet structures
(\texttt{TelemetryPacket}, \texttt{BeaconPacket}, etc.) are defined
once in \texttt{tdma\_protocol.h} with \texttt{\#pragma pack(1)}.

A hardware abstraction layer (HAL) in \texttt{src/hal/} wraps all
platform-exclusive APIs (reboot, RNG, NVS, AES, task creation,
spinlock) behind thin C~function headers.\ Application code above the
HAL is platform-agnostic; only the four \texttt{hal/*.cpp}
implementation files differ between the ESP32 and STM32 targets.

%% ============================================================
\section{Validation Plan}
\label{sec:validation}

\subsection{Phase 1: Analog Front-End Characterisation}

The analog front-end will be characterised against a calibrated
Fluke~435-II power quality analyser as the reference.\ Test
conditions cover:
\begin{itemize}
  \item Voltage accuracy: mains voltage varied $\pm 10\%$ around 220~V
        via a variac; target \(\le 0.5\%\) error.
  \item Current accuracy: precision resistive loads (100~W bulb,
        1000~W heater) measured at PF~=~1.0; target \(\le 0.5\%\) error.
  \item Power factor accuracy: ceiling fan and window-type aircon
        (PF~=~0.55--0.90); target \(\le 0.005\).
  \item Phase calibration: empirical determination of CT phase error
        at 5~A and 30~A load levels using the two-point PHASE register
        correction procedure.
\end{itemize}

\subsection{Phase 2: Load Classification Validation}

The decision tree classifier will be validated against 18 physical
loads (3 examples per class, listed in Table~\ref{tab:testloads}).
Acceptance criteria: \(\ge 85\%\) class-averaged accuracy; no single
class below 70\%.\ After validation, supervised labelling data will be
collected (50--100 feature vectors per load) and used to train the
TFLite Micro MLP offline.\ The trained model will be deployed via the
GFSK bulk channel and re-evaluated against the same test set.

\begin{table}[t]
  \caption{Load Classification Validation Test Set}
  \label{tab:testloads}
  \centering
  \small
  \begin{tabular}{lp{5.2cm}}
    \toprule
    Class & Physical Test Loads \\
    \midrule
    RESISTIVE  & 100~W incandescent, electric flat iron, space heater \\
    CAPACITIVE & Active-PFC laptop charger, UPS (battery mode) \\
    MOTOR\_1PH & 0.75~HP window aircon, refrigerator compressor,
                  washing machine \\
    FAN        & 75~W ceiling fan, desk fan, exhaust fan \\
    SMPS       & Laptop charger (no PFC), desktop PC PSU, 5~W USB charger \\
    LIGHTING   & 13~W CFL, 9~W LED bulb (non-PFC),
                  36~W fluorescent + magnetic ballast \\
    \bottomrule
  \end{tabular}
\end{table}

\subsection{Phase 3: Network Performance Characterisation}

With all 8 node slots active and transmitting simulated telemetry
(\texttt{PZEM\_FAKE} build flag), the uplink packet delivery ratio
(PDR) will be measured over 1000 consecutive superframes (50~min) in
a residential indoor environment (concrete walls, kitchen/living room
separation) at typical node-to-gateway distances of 5--20~m.\ Target:
PDR \(\ge 99\%\).\ Latency from sensor sampling to gateway WebSocket
push will be measured; target: \(<\) 3.5~s (one superframe cycle).

\subsection{Phase 4: GFSK Bulk Transfer Validation}

The GFSK bulk transport will be validated through:
\begin{enumerate}
  \item Echo test (DL + UL round-trip) at payload sizes 1~B, 200~B,
        201~B, 500~B, and 4096~B.
  \item Forced-retransmission test: deliberate NACK injection to verify
        ARQ recovery and end-to-end CRC-32 integrity.
  \item Overrun protection: verify the GFSK phase terminates before
        \texttt{END\_GUARD\_MS} under maximum payload + maximum retries.
  \item Telemetry PDR regression: PDR must remain \(\ge 99\%\) during
        bulk-active superframes.
\end{enumerate}

\subsection{Phase 5: MLP Deployment and Re-Evaluation}

The quantised TFLite Micro MLP model (\(\approx 500\)~bytes) will be
deployed to all 8 nodes via the GFSK broadcast mechanism
(\texttt{POST /api/bulk/broadcast}).\ Classification accuracy will be
re-measured using the same 18-load test set.\ Target: \(\ge 90\%\)
class-averaged accuracy with the MLP, compared to the decision tree
baseline.

%% ============================================================
\section{Expected Contributions and Outcomes}
\label{sec:contributions}

The proposed system makes the following contributions:

\begin{enumerate}
  \item \textbf{Precision waveform-accessible metering:} The
        ADS131M02 + ZMPT101B + CT front-end provides simultaneous
        voltage and current sampling at 4~kSPS with \(\le 0.5\%\)
        accuracy and full galvanic isolation, enabling harmonic analysis
        that is impossible with the PZEM-004T.

  \item \textbf{Scalable deterministic MAC:} The frequency-hopping
        TDMA protocol supports up to 8 nodes in a 3~s superframe with
        guaranteed slot timing, collision-free operation, and a
        provably unique per-packet keystream for AES-128 CTR
        encryption.

  \item \textbf{On-node, no-transmission-cost load classification:}
        The decision tree and MLP classifiers run entirely on the
        STM32F411 node; only a 1-byte result is transmitted per
        superframe.\ This avoids the latency and bandwidth cost of
        raw waveform uplink and maintains classification capability
        even when the gateway is unreachable.

  \item \textbf{OTA model management:} The GFSK bulk transfer
        extension enables deployment of updated classifier models and
        calibration data to all nodes within 24~s without disrupting
        ongoing telemetry, providing a practical path from
        decision-tree to trained MLP operation.

  \item \textbf{Self-hosted residential dashboard:} The web SPA
        requires no cloud infrastructure; the gateway operates as an
        access point or connects to an existing router, making the
        system deployable in locations without internet connectivity.
\end{enumerate}

The system is intended to serve as a research platform for evaluating
harmonic-signature load classification under Philippine residential
conditions (220~V / 60~Hz, tropical climate, diverse load mix).\ Data
collected through the supervised labelling mode will contribute to a
small labelled power signature dataset for the Philippine residential
context, where public datasets such as REDD \cite{kolter_redd_2011}
and UK-DALE \cite{kelly_neural_nilm_2015} are not directly applicable
due to differing mains frequency (60~Hz vs.\ 50~Hz) and appliance
profiles.

%% ============================================================
\section{Conclusion}
\label{sec:conclusion}

This paper has proposed a distributed power telemetry and load
classification system tailored for Philippine residential energy
monitoring.\ The system replaces the commodity PZEM-004T module with
a precision ADS131M02 analog front-end that provides simultaneous
voltage/current sampling, full galvanic isolation, and PGA
gain-ranging across a 5000:1 current dynamic range.\ A TDMA LoRa MAC
with frequency hopping, AES-128 CTR encryption, and a novel GFSK
bulk transfer extension provides scalable, secure wireless
connectivity for up to 8 nodes.\ An on-node harmonic feature
extraction pipeline feeding a deterministic decision tree (with a
TFLite Micro MLP upgrade path) classifies connected loads into six
residential categories without transmitting raw waveform data.

The integrated system is designed to validate \(\le 0.5\%\) metering
accuracy, \(\ge 85\%\) classification accuracy, and \(\ge 99\%\)
packet delivery ratio under realistic residential conditions.\ The
proposed work addresses gaps in existing low-cost metering and NILM
literature by combining precision simultaneous-channel ADC metering,
structured encrypted wireless MAC, and embedded load classification in
a single self-hosted platform applicable to the 60~Hz Philippine
residential context.

%% ============================================================
\begin{thebibliography}{00}

\bibitem{hart_nilm_1992}
G.\ W.\ Hart, ``Nonintrusive appliance load monitoring,'' \textit{Proc.\
IEEE}, vol.~80, no.~12, pp.~1870--1891, Dec.\ 1992.

\bibitem{zeifman_nilm_review_2011}
M.\ Zeifman and K.\ Roth, ``Nonintrusive appliance load monitoring:
Review and outlook,'' \textit{IEEE Trans.\ Consum.\ Electron.}, vol.~57,
no.~1, pp.~76--84, Feb.\ 2011.

\bibitem{wichakool_harmonic_loads_2009}
W.\ Wichakool, A.-T.\ Avestruz, R.\ W.\ Cox, and S.\ B.\ Leeb,
``Modeling and estimating current harmonics of variable electronic
loads,'' \textit{IEEE Trans.\ Power Electron.}, vol.~24, no.~12,
pp.~2803--2811, Dec.\ 2009.

\bibitem{lam_vI_trajectory_2007}
H.\ Y.\ Lam, G.\ S.\ K.\ Fung, and W.\ K.\ Lee, ``A novel method to
construct taxonomy of electrical appliances based on load signatures,''
\textit{IEEE Trans.\ Consum.\ Electron.}, vol.~53, no.~2,
pp.~653--660, May 2007.

\bibitem{kolter_redd_2011}
Z.\ Kolter and M.\ Johnson, ``REDD: A public data set for energy
disaggregation research,'' in \textit{Proc.\ Workshop Data Mining Appl.\
Sustain.\ (SustKDD)}, 2011.

\bibitem{kelly_neural_nilm_2015}
J.\ Kelly and W.\ Knottenbelt, ``Neural NILM: Deep neural networks
applied to energy disaggregation,'' in \textit{Proc.\ 2nd ACM Int.\ Conf.\
Embedded Syst.\ Energy-Efficient Built Environ.\ (BuildSys)}, 2015,
pp.~55--64.

\bibitem{parson_unsupervised_nilm_2012}
O.\ Parson, S.\ Ghosh, M.\ Weal, and A.\ Rogers, ``An unsupervised
training method for non-intrusive appliance load monitoring,''
\textit{Artif.\ Intell.}, vol.~217, pp.~1--19, Dec.\ 2014.

\bibitem{adamo_power_quality_2011}
F.\ Adamo, F.\ Attivissimo, A.\ Di Nisio, and M.\ Spadavecchia,
``A low-cost tool for power quality monitoring in home automation
systems,'' in \textit{Proc.\ IEEE Instrum.\ Meas.\ Technol.\ Conf.\
(I2MTC)}, 2011, pp.~1--6.

\bibitem{ti_ads131m02_ds}
Texas Instruments, ``ADS131M02: 2-channel, simultaneously sampling,
24-bit, delta-sigma ADC,'' Datasheet SBAS865B, Sep.\ 2020.\ [Online].
Available: \url{https://www.ti.com/product/ADS131M02}

\bibitem{zmpt101b_datasheet}
Zhongce Transformer Co., ``ZMPT101B precision miniature voltage
transformer,'' Component Datasheet, 2018.

\bibitem{augustin_lora_sensors_2016}
A.\ Augustin, J.\ Yi, T.\ Clausen, and W.\ M.\ Townsley, ``A study of
LoRa: Long range and low power networks for the internet of things,''
\textit{Sensors}, vol.~16, no.~9, p.~1466, Sep.\ 2016.

\bibitem{raza_lpwan_survey_2017}
U.\ Raza, P.\ Kulkarni, and M.\ Sooriyabandara, ``Low power wide area
networks: An overview,'' \textit{IEEE Commun.\ Surv.\ Tut.}, vol.~19,
no.~2, pp.~855--873, 2nd Qtr.\ 2017.

\bibitem{bor_lorawan_scale_2016}
M.\ Bor, J.\ Vidler, and U.\ Roedig, ``LoRa for the internet of
things,'' in \textit{Proc.\ 13th Int.\ Workshop Mobile Ubiquitous Syst.:
Comput., Netw.\ Services (EWsN)}, 2016.

\bibitem{semtech_sx1278_ds}
Semtech Corporation, ``SX1276/77/78/79: 137 MHz to 1020 MHz low power
long range transceiver,'' Datasheet Rev.~7, 2020.\ [Online].
Available: \url{https://www.semtech.com/products/wireless-rf/lora-core/sx1276}

\bibitem{nist_sp800_38a}
M.\ Dworkin, ``Recommendation for block cipher modes of operation:
Methods and techniques,'' NIST Special Publication 800-38A, Dec.\
2001.\ [Online].\ Available: \url{https://doi.org/10.6028/NIST.SP.800-38A}

\bibitem{lorawan_spec}
LoRa Alliance, ``LoRaWAN\texttrademark{} 1.0.4 specification,'' Oct.\
2020.\ [Online].\ Available: \url{https://lora-alliance.org/resource_hub/lorawan-104-specification-package/}

\bibitem{hancke_rfid_security_2005}
G.\ P.\ Hancke and M.\ G.\ Kuhn, ``An RFID distance bounding protocol,''
in \textit{Proc.\ 1st Int.\ Conf.\ Security Privacy Emerging Areas
Commun.\ Netw.\ (SecureComm)}, 2005, pp.~67--73.

\bibitem{david_tflite_micro_2021}
R.\ David \textit{et al.}, ``TensorFlow Lite Micro: Embedded machine
learning for TinyML systems,'' \textit{Proc.\ Mach.\ Learn.\ Syst.
(MLSys)}, 2021.\ [Online].\ Available: \url{https://arxiv.org/abs/2010.08678}

\bibitem{warden_tinyml_2019}
P.\ Warden and D.\ Situnayake, \textit{TinyML: Machine Learning with
TensorFlow Lite on Arduino and Ultra-Low-Power Microcontrollers}.\
Sebastopol, CA, USA: O'Reilly Media, 2019.

\bibitem{arm_cmsis_dsp}
ARM Limited, ``CMSIS-DSP software library,'' Documentation Version
1.15.0, 2024.\ [Online].\ Available:
\url{https://arm-software.github.io/CMSIS_5/DSP/html/index.html}

\bibitem{radiolib}
J.\ Gromes, ``RadioLib: Universal wireless communication library for
Arduino,'' GitHub, 2024.\ [Online].\ Available:
\url{https://github.com/jgromes/RadioLib}

\bibitem{ieee_std_1459}
IEEE, ``IEEE standard definitions for the measurement of electric power
quantities under sinusoidal, nonsinusoidal, balanced, or unbalanced
conditions,'' IEEE Std 1459-2010, Feb.\ 2010.

\bibitem{stm32f411_rm}
STMicroelectronics, ``STM32F411xC/xE reference manual,'' RM0383
Rev.~4, 2023.\ [Online].\ Available:
\url{https://www.st.com/resource/en/reference_manual/rm0383-stm32f411xce-advanced-armbased-32bit-mcus-stmicroelectronics.pdf}

\end{thebibliography}

\end{document}
